\textbf{(CAP-COLUNI 2020 -- Q19)} Considere as afirmativas abaixo, atribuindo V para a(s) verdadeira(s) e F para a(s) falsa(s):

\begin{itemize}
\item[(\ \ )] As raízes da equação do 2º grau $x^2-4x+3=0$ são números primos.
\item[(\ \ )] Dados: $\mathbb{Z}$ -- conjunto dos números inteiros, $\mathbb{Q}$ -- conjunto dos números racionais e $\mathbb{R}$ -- conjunto dos números reais, pode-se afirmar que $C_{\mathbb{R}}^{\mathbb{Z}}=\mathbb{Q}$, sendo $C_{\mathbb{R}}^{\mathbb{Z}}=\mathbb{Q}$ o complementar de $\mathbb{Z}$ em $\mathbb{R}$.
\item[(\ \ )] O produto das raízes da equação irracional
\[
\sqrt{2x^3-3x^2+6x}=x\sqrt{2x+7}
\]
é zero.
\item[(\ \ )] Sendo $x=\sqrt{2+\sqrt3}$ e $y=\sqrt{2-\sqrt3}$, o produto $x\cdot y$ é igual a $\pm1$.
\end{itemize}

Assinale a sequência CORRETA:

\begin{enumerate}
	\item F, F, V, F.
	\item F, V, F, V.
	\item V, F, V, F.
	\item V, V, F, V.
\end{enumerate}
